\section{Assembly and the fixed-time crossed limit}
\label{sec:assembly}
This section completes the compact distribution step for the specified
equal-species, strict-$G_N\to0$ model. We use the exact four terms in
KL (B.21)--(B.24), rather than dropping the second endpoint branch of
an individual hypergeometric function. The delta coefficient agrees with
KL (B.46); the purpose here is to control the combined remainder and the
finite-packet limit. No originality claim or finite-$G_N$ result is intended.

\subsection{Finite parameters and the four coefficients}
Put $s_j=\Lambda+q_j$, with real offsets in a fixed compact set, and define
\begin{align}
 p&=\tfrac12+\tfrac i2(q_4-q_2-2\nu),&
 q&=\tfrac12+\tfrac i2(q_4-q_2+2\nu),\notag\\
 e&=i(q_1-q_3),& \delta&=i(q_1-q_2+q_3-q_4),\notag\\
 t&=(\delta-e)/2,&v&=(\delta+e)/2.
 \label{eq:assembly-vars}
\end{align}
Here the scalar $q$ is distinguished from the four offsets $q_j$.
Write
\begin{equation}
 c(z)=\sqrt\pi\,\frac{\Gamma((1-z)/2)}{\Gamma(z/2)}.
 \qquad c(z)c(1-z)=\pi.
 \label{eq:assembly-c}
\end{equation}
The beta-like function of KL (B.23) is exactly
$B(u,w)=c(u)c(w)/c(u+w-1)$. Uniform gamma-ratio estimates give
\begin{equation}
 B(u,\pm2i\Lambda+w_0)
 =c(u)\Lambda^{u-1}[1+O(\Lambda^{-1})]
 \label{eq:assembly-beta-asym}
\end{equation}
on compact offset sets away from the indicated meromorphic poles.
There is no arbitrary constant in the phase of this formula: the powers
come from the principal logarithms of the two opposite large arguments.

Set $D=c(1-p-t)c(q+v)c(1-q+t)c(p-v)$. Denote the exact prefactors of
the four hypergeometric functions by $P_{j,\Lambda}$, in the order
of the two terms of the first $G$, then the two terms of the second $G$
in KL (B.21). Thus $\pi^3S=\sum_jP_{j,\Lambda}H_{j,\Lambda}$.
The complete parameter and leading-prefactor table is
\begin{center}
\small
\begin{tabular}{c c c c c c}
\toprule
$j$ & $b_j$ & $d_j$ & $g_j$ & $\delta_j$ & $W_j$\\
\midrule
1 & $1-p-v$ & $p-v$ & $1-e$ & $\delta$ & $-2i\Lambda$\\
2 & $1-q+t$ & $q+t$ & $1-e$ & $-\delta$ & $2i\Lambda$\\
3 & $1-q+v$ & $q+v$ & $1+e$ & $-\delta$ & $2i\Lambda$\\
4 & $1-p-t$ & $p-t$ & $1+e$ & $\delta$ & $-2i\Lambda$\\
\bottomrule
\end{tabular}
\end{center}
\begin{align}
 P_{j,\Lambda}&=\Lambda^{\delta_j-1}P_j^0[1+O(\Lambda^{-1})],\notag\\
 (P_1^0,P_2^0,P_3^0,P_4^0)
 &=D\left(\frac{c(p+v)}{c(p-v)c(e)},
 \frac{c(1-q-t)}{c(1-q+t)c(e)},
 \frac{c(q-v)}{c(q+v)c(-e)},
 \frac{c(1-p+t)}{c(1-p-t)c(-e)}\right).
 \label{eq:assembly-prefactors}
\end{align}
For example, write the original $h$ variables as
$h_1=p$, $h_3=q$, $h_2=H$, $h_4=1-H$, $h=K$, $h'=L$.
Then $H-K=t$, $H-L=v$, $K-L=e$. The first prefactor consists of
$B(1+H-K-q,2-2H)$, $B(1+K-p-H,p+H-L)$,
and $B(H+q-L,K+H-q)$. Applying \eqref{eq:assembly-beta-asym}
to the first and third yields $P_1^0$. The second term follows by
the other choice of $h_4$, and the second $G$ by
$(p,q,H,K,L)\mapsto(q,p,1-H,1-K,1-L)$.
This determines the table without choosing branches by numerical size.

\subsection{Cancellation and a pole-free leading expression}
The additional endpoint branch of term $j$, after multiplying by its
leading prefactor, has coefficient
\begin{equation}
 E_j=-\frac{P_j^0\Gamma(g_j)}
 {\delta_j\Gamma(b_j)\Gamma(d_j)}
 (W_j/\Lambda)^{-\delta_j}
 \label{eq:assembly-extra}
\end{equation}
in $\Lambda\pi^3S$. The identity $\sum_jE_j=0$ is exact as a meromorphic
identity. To check it explicitly, put
$P=\pi p/2$, $Q=\pi q/2$, $T=\pi t/2$, $V=\pi v/2$.
Using $c(z)=2^z\Gamma(1-z)\sin(\pi z/2)$, their sum is
$2D e^{i\pi\delta/2}/[\pi\delta\sin(\pi e/2)]$ times
\begin{align}
 &-\sin(P+V)\cos(P-V)+\cos(Q+T)\sin(Q-T)\notag\\
 &\hspace{8mm}-\sin(Q-V)\cos(Q+V)+\cos(P-T)\sin(P+T)=0.
 \label{eq:assembly-trig-zero}
\end{align}
The equality follows by writing each product as a sum of two sines.

The remaining coefficients also simplify. Define
\begin{align}
 U_+&=2^{1-\delta}\pi\cos(\pi\delta/2)
           \frac{c(p+v)c(q+v)}{c(p-t)c(q-t)},\notag\\
 U_-&=2^{1+\delta}\pi\cos(\pi\delta/2)
           \frac{c(p-v)c(q-v)}{c(p+t)c(q+t)}.
 \label{eq:assembly-U}
\end{align}
They equal, respectively,
$\sum_{j=1,4}P_j^0\Gamma(g_j)/[\Gamma(g_j-b_j)\Gamma(g_j-d_j)]$
and the analogous sum over $j=2,3$. One way to verify the first identity
is to factor
$D/[c(p-v)c(p-t)\Gamma(p+t)\Gamma(p+v)]$ from the two terms.
The remaining numerator is
\[
 \sin(P-T)\cos(P+T)-\sin(P-V)\cos(P+V)
 =\cos(\pi\delta/2)\sin(\pi e/2),
\]
which cancels the apparent $e$ pole. The second identity follows by the
same interchange as above. In particular,
\begin{equation}
 U_+|_{\delta=0}=U_-|_{\delta=0}=2\pi
 \label{eq:assembly-residue}
\end{equation}
for every allowed $p,q,e$, including $e=0$.

\begin{theorem}[Uniform assembled remainder]
On compact real-offset sets, with fixed $\nu>0$, for any
$0<\epsilon<1/4$,
\begin{align}
 S(\Lambda+q_1,\ldots,\Lambda+q_4,\nu,\nu)
 &=S_\Lambda^0+O(\Lambda^{-2+\epsilon}),\notag\\
 S_\Lambda^0&=\frac{1}{\pi^3\Lambda}
   [\Lambda^\delta\Gamma(\delta)U_+
       +\Lambda^{-\delta}\Gamma(-\delta)U_-].
 \label{eq:assembled-asymptotic}
\end{align}
The bracket has its removable value at $\delta=0$; the estimate includes
both opposite-label coincidence surfaces and their intersection.
\label{thm:assembled-remainder}
\end{theorem}
\begin{proof}
First stay away from $e=0$ and $\delta=0$. Multiplying the endpoint
expansion of Proposition~\ref{prop:endpoint} by
\eqref{eq:assembly-prefactors}, and using the exact cancellation above,
gives the displayed remainder. To pass through the apparent poles,
we require estimates in a small complex tube, not just continuity on
the real labels.

The Euler/Fourier proof extends to such tubes: the real parts of
$A_0,C_0,\alpha,\beta,b,d,g-d$ remain positive, the balanced identity
$\alpha+\beta+\delta=1$ remains exact, and the triple-corner integrability
is unchanged under sufficiently small perturbations. On compact subsets
of this tube separated from $\delta=0$, that proof gives a remainder
$O(\Lambda^{-1+\eta})$ when $|\Re\delta|\le\eta$.
The prefactors have size $O(\Lambda^{-1+\eta})$ off $e=0$, while their
relative gamma-ratio error is $O(\Lambda^{-1})$. Thus the combined
remainder is $O(\Lambda^{-2+2\eta})$ there. Choose $2\eta<\epsilon$.

The exact $S$ is holomorphic in a mass-independent small tube about any
compact real-offset set, by its real triple integral KL (B.10)--(B.14).
For completeness, use signed logarithmic coordinates for its three
integration variables. The real-label envelope is an exponential
decay times the four hyperplane singularities
$|u|^{-1/2}$, $|v|^{-1/2}$, $|w|^{-1/2}$, $|u+v+w|^{-1/2}$.
A small imaginary perturbation of the labels weakens the exponential
rate and changes each singular exponent by at most $C\eta$.
Choose the tube so that each exponent is below one and their sum at
the common origin is below three, while the exponential rate stays
positive. The same envelope with these margins dominates parameter
derivatives, proving holomorphy. The large real mass appears only in
oscillatory phases and does not shrink this tube.

The explicit $U_\pm$ are holomorphic throughout the tube: their $c$
arguments remain near real part $1/2$. Equation~\eqref{eq:assembly-residue}
removes the $\delta$ pole, and there is no $e$ denominator. Hence
$S-S_\Lambda^0$ is holomorphic in $e$ and $\delta$ across the relevant
surfaces. Apply the maximum-modulus principle on small disks in these
variables, using a product of circles near their intersection. The
boundary estimates just obtained transfer to the interiors. Compact
covering gives the asserted real-parameter estimate. In this chart
only $e=0$ is an apparent prefactor pole; the other opposite-label
surface introduces none, so is already included.
\end{proof}

For numerical evaluation at $\delta=0$, define
$\ell(z)=\partial_z\log c(z)
=-[\psi((1-z)/2)+\psi(z/2)]/2$. Then
\begin{equation}
 S_\Lambda^0\big|_{\delta=0}
 =\frac{4}{\pi^2\Lambda}
 \left[\log(\Lambda/2)-\gamma_E
 +\frac12\sum_{z=p,q}\{\ell(z+e/2)+\ell(z-e/2)\}\right].
 \label{eq:assembly-zero-value}
\end{equation}
No coincident-label perturbation is needed in this formula.

\subsection{The local distribution and exact-energy variables}
Write $\chi=q_1-q_2+q_3-q_4$, so $\delta=i\chi$.
Equation~\eqref{eq:assembly-residue} isolates the singular part of
$\Lambda S_\Lambda^0$ as
\begin{equation}
 \frac{4}{\pi^2}\frac{\sin(\chi\log\Lambda)}{\chi}.
 \label{eq:assembly-sinc}
\end{equation}
After subtracting this expression, the two remaining amplitudes are
smooth through $\chi=0$ and multiply $e^{\pm i\chi\log\Lambda}$.
Their pairing with a compact smooth test tends to zero by integration
by parts. The sinc distribution therefore proves
\begin{equation}
 \Lambda S\longrightarrow\frac4\pi\delta(\chi)
 \quad\hbox{in distributions on the four offsets}.
 \label{eq:raw-S-delta}
\end{equation}
This also shows $\Lambda S=O(1+\log\Lambda)$ pointwise on compact real
offset sets, a useful bound when changing coordinates.

Return to the ordering $(a,b,c,d)$ used in our conditional kernel and
write $\mathcal C_\nu=F_\Lambda S$. Directly from A.14 and
Section~\ref{sec:fourpoint},
\begin{equation}
 F_\Lambda=
 \frac{abc\sinh(\pi a)\tanh(\pi b)\sinh(\pi c)\cosh^2(\pi d)}
 {1024\pi^9}\,\Gamma_{16}(a,b,c,d).
 \label{eq:assembly-external-F}
\end{equation}
Here $\Gamma_{16}$ is exactly the sixteen-gamma product of A.14.
Its eight large arguments give
$16\pi^4\Lambda^{-2}e^{-\pi(a+b+c+d)}[1+O(\Lambda^{-1})]$;
the remaining eight factors, denoted $\Gamma_{\rm diff}$, depend only
on adjacent label differences. Consequently
\begin{equation}
 \frac{F_\Lambda}{\Lambda}
 \longrightarrow\frac{e^{\pi(d-b)}}{1024\pi^5}\Gamma_{\rm diff}
 \label{eq:assembly-F-limit}
\end{equation}
smoothly on compact offsets, with an $O(\Lambda^{-1})$ remainder in
every fixed derivative order. The phases $\Lambda^{i\eta\nu}$ of the
four large gamma pairs cancel because $\sum\eta=0$.

On $a=d+x$, $c=d+z$, $b=d+x+z$, multiplication of
\eqref{eq:raw-S-delta} by \eqref{eq:assembly-F-limit} gives
\begin{equation}
 \frac{e^{-\pi(x+z)}}{256\pi^6}
 \prod_{u=x,z}\left|
 \Gamma\left(\tfrac14+\tfrac i2(\nu+u)\right)
 \Gamma\left(\tfrac14+\tfrac i2(\nu-u)\right)\right|^2
 =v_\nu(x)v_\nu(z).
 \label{eq:assembly-delta-weight}
\end{equation}
The two orderings of $\chi$ differ by a sign only, leaving the delta
measure unchanged.

Finally, the exact-energy change $s(M+\xi)-M=\xi+O(M^{-1})$ is a
smooth diffeomorphism close to the identity on every compact set.
In a four-variable test pairing, its Jacobian cancels the three internal
factors in \eqref{eq:translated-cross-density}, leaving only the inverse
external Jacobian, which tends smoothly to one. The transformed test
differs by $O(M^{-1})$ in every fixed smooth seminorm. The compact
$O(1+\log M)$ bound just proved makes this difference vanish even
against the moving kernel. Thus we obtain the formerly conditional step:
\begin{theorem}[Compact limit and finite-packet consequence]
For the model and normalizations of Sections~\ref{sec:fourpoint} and
\ref{sec:crosslimit}, equation~\eqref{eq:local-racah-hypothesis} holds.
Proposition~\ref{prop:conditional-cross-limit} therefore applies without
its extra hypothesis: both threshold protocols have their stated
fixed-time limits, uniformly on compact sets of real times, for every
fixed probability density $\pi_0\in L^1(\mathbb R_+)$.
\label{thm:fixed-time-crossed}
\end{theorem}
The moving-threshold commutator limit is the positive static term
\eqref{eq:comparison-boundary-floor}: the ordered norms converge by
the one-step gamma-ratio limit and the same Gaussian domination.
For a fixed threshold the repeated-time commutator limit is zero.
These statements retain the fixed Gaussian width and strict $G_N\to0$.
They do not determine the scrambling-time remainder, a Lyapunov exponent,
or a finite-$G_N$ holographic dual.
